\documentclass[11pt]{article}

\usepackage[margin=1in]{geometry}
\usepackage{amsmath,amssymb,amsthm,mathtools}
\usepackage{microtype}
\usepackage[colorlinks=true,linkcolor=blue,citecolor=blue,urlcolor=blue]{hyperref}

\newtheorem{theorem}{Theorem}[section]
\newtheorem{proposition}[theorem]{Proposition}
\newtheorem{lemma}[theorem]{Lemma}
\newtheorem{remark}[theorem]{Remark}

\newcommand{\C}{\mathbb C}
\newcommand{\D}{\mathbb D}
\newcommand{\Mn}{\mathbb M_n(\C)}
\newcommand{\RePart}{\operatorname{Re}}
\newcommand{\alg}{\operatorname{alg}}
\newcommand{\diag}{\operatorname{diag}}
\newcommand{\T}{\mathsf T}

\title{Structured Positive-Kernel Elimination\\
and Gramian Rigidity as a Reusable Toolkit}
\author{Research note}
\date{}

\begin{document}

\maketitle

\begin{abstract}
Assuming that the candidate proof of the scalar Crouzeix conjecture
survives expert review, its central mechanism can be separated into a
reusable two-stage method: structured elimination of nuisance terms
from a positive kernel, followed by a rigidity argument for two
weighted Gramians.  The method is not valid for arbitrary positive
kernels.  Its essential hypothesis is an algebraic alignment that
makes the nuisance module exactly annihilable.
\end{abstract}

\section{The guiding abstraction}

The proposed toolkit may be summarized as
\begin{equation}
 \boxed{\text{structured positive-kernel elimination}
 \quad\Longrightarrow\quad
 \text{Gramian rigidity}.}
 \label{eq:guiding-abstraction}
\end{equation}
The first stage removes an unknown but structured term without an
estimate.  The second stage converts the surviving positive-matrix
inequality into a sharp norm bound through detectability and a Stein
identity.

\section{A kernel-elimination principle}

Suppose that
\begin{equation}
 H(z)=F(z)+E(z),
 \qquad \RePart H(z)\succeq0,
 \label{eq:H-splitting}
\end{equation}
where $F$ is known and the nuisance term $E$ is constrained to a
linear space \(\mathcal E\).  Positivity of the real part produces the
Herglotz kernel
\begin{equation}
 L_H(z,w)=\frac{H(z)+H(w)^*}{1-z\overline w}\succeq0.
 \label{eq:Herglotz-kernel}
\end{equation}

Choose points $z_1,\ldots,z_N\in\D$ and vectors
\(\xi_1,\ldots,\xi_N\), and set
\begin{equation}
 \eta_i=\sum_{j=1}^N
 \frac{\xi_j}{1-z_i\overline{z_j}}.
 \label{eq:eta-definition}
\end{equation}
The contribution of $E$ to the sampled kernel quadratic form is
exactly
\begin{equation}
 2\RePart\sum_{i=1}^N \xi_i^*E(z_i)\eta_i.
 \label{eq:nuisance-contribution}
\end{equation}
Consequently, test data satisfying
\begin{equation}
 \RePart\sum_{i=1}^N \xi_i^*E(z_i)\eta_i=0
 \qquad(E\in\mathcal E)
 \label{eq:annihilator-condition}
\end{equation}
remove the entire nuisance module from the positive-kernel
inequality.  What remains depends only on the known part $F$.

In the candidate Crouzeix proof, after an adjoint similarity the
nuisance module has the form
\begin{equation}
 E(z)=D(z)G,
 \qquad D(z)\ \text{diagonal},
 \label{eq:diagonal-nuisance}
\end{equation}
and the auxiliary vector $v$ is determined by the single equation
\begin{equation}
 Gv+Pu=0.
 \label{eq:annihilator-equation}
\end{equation}
This equation constructs an annihilator for every diagonal analytic
correction simultaneously.  The step may therefore be regarded as an
\emph{annihilator compression} of a positive kernel.

\section{A standalone Gramian-rigidity lemma}

The second part of the mechanism is independent of numerical ranges.

\begin{lemma}[Detector--anticommutator rigidity]
\label{lem:detector-anticommutator}
Let $P=P^*$ and $X=X^*$ be finite matrices satisfying
\begin{equation}
 P\succeq I,
 \qquad X\succeq0,
 \qquad cX-XP-PX\succeq0,
 \qquad c\geq2.
 \label{eq:anticommutator-hypotheses}
\end{equation}
Assume the detectability condition
\begin{equation}
 \ker X\subseteq\ker(P-I).
 \label{eq:detectability}
\end{equation}
Then
\begin{equation}
 P\preceq\frac c2 I.
 \label{eq:P-rigidity-bound}
\end{equation}
If, in addition,
\begin{equation}
 P=I+aC^*PC,
 \qquad a>0,
 \label{eq:Stein-equation}
\end{equation}
then
\begin{equation}
 \|C\|^2\leq\frac{c/2-1}{a}.
 \label{eq:C-bound}
\end{equation}
\end{lemma}

\begin{proof}
Let $Pe=pe$ for a unit eigenvector $e$.  Testing the third
inequality in \eqref{eq:anticommutator-hypotheses} on $e$ gives
\begin{equation}
 0\leq(c-2p)e^*Xe.
 \label{eq:eigenvector-test}
\end{equation}
If $p>c/2$, positivity of $X$ forces $e^*Xe=0$, and hence
$Xe=0$.  Detectability then gives $Pe=e$, contradicting
$p>c/2\geq1$.  Thus every eigenvalue of $P$ is at most $c/2$,
which proves \eqref{eq:P-rigidity-bound}.

Finally, $P\succeq I$ and \eqref{eq:Stein-equation} imply
\[
 aC^*C\preceq aC^*PC=P-I
 \preceq\left(\frac c2-1\right)I.
\]
This is \eqref{eq:C-bound}.
\end{proof}

The same detector argument extends to a completely positive dynamic
\begin{equation}
 \Phi(Y)=\sum_j C_j^*YC_j.
 \label{eq:CP-dynamic}
\end{equation}
Whenever the corresponding Stein equation and detectability
condition hold, it gives bounds on the row energy
\(\sum_jC_j^*C_j\).  This makes
Lemma~\ref{lem:detector-anticommutator} potentially useful in
multivariable realization and control problems.

\section{Scale selection and the distinguished radius
\texorpdfstring{$1/2$}{1/2}}

The two scales in the candidate proof arise from an optimization
principle.  For $0<r<1$, sampling at
\begin{equation}
 z_i=r\overline{\lambda_i}
 \label{eq:r-sampling}
\end{equation}
produces the weighted Gramians
\begin{align}
 P_r&=\sum_{k\geq0}r^{2k}C^{*k}C^k,
 \label{eq:Pr}\\
 R_r&=\sum_{k\geq0}r^kC^{*k}C^k,
 \label{eq:Rr}\\
 X_r&=R_r-P_r
 =\sum_{k\geq1}(r^k-r^{2k})C^{*k}C^k\succeq0.
 \label{eq:Xr}
\end{align}
Exact cancellation yields
\begin{equation}
 \frac{2}{1-r}X_r-X_rP_r-P_rX_r\succeq0.
 \label{eq:r-anticommutator}
\end{equation}
The detector condition is automatic here: $X_re=0$ implies
$Ce=0$, and then the series for $P_r$ gives $P_re=e$.
Lemma~\ref{lem:detector-anticommutator}, with
\(c=2/(1-r)\), therefore gives
\begin{equation}
 P_r\preceq\frac1{1-r}I.
 \label{eq:Pr-bound}
\end{equation}
The Stein identity
\begin{equation}
 P_r=I+r^2C^*P_rC
 \label{eq:Pr-Stein}
\end{equation}
then yields
\begin{equation}
 \|C\|\leq\frac1{\sqrt{r(1-r)}}.
 \label{eq:r-norm-bound}
\end{equation}
Since $r(1-r)$ is uniquely maximized at $r=1/2$, the optimal
choice is
\begin{equation}
 r=\frac12,
 \qquad \|C\|\leq2.
 \label{eq:optimal-radius}
\end{equation}
Thus $1/2$ is not merely a fortunate sampling value: it is the
unique optimizer of the two-scale family.

\section{Immediate structural extensions}

The simple-spectrum assumption is not the essential feature of the
cancellation.  Suppose that
\begin{equation}
 T=S\Lambda S^{-1}
 \label{eq:diagonalization}
\end{equation}
is diagonalizable, possibly with repeated eigenvalues, and that the
correction belongs after adjoint similarity to a row-separable
diagonal module
\begin{equation}
 \mathcal E_S=
 \left\{S^{-*}DS^*:D\ \text{diagonal}\right\}.
 \label{eq:row-separable-module}
\end{equation}
The same rowwise annihilator construction applies.  If the correction
belongs specifically to \(\alg(T^*)\), repeated eigenvalues merely
impose additional equalities among the relevant diagonal entries.

This observation points toward extensions to finite-dimensional
semisimple commutative algebras and simultaneously diagonalizable
commuting tuples.  Removing diagonalizability would require an
additional argument, such as radial regularization followed by a
controlled perturbation; this is a proposed direction rather than a
consequence established here.

\section{Promising open-problem targets}

\subsection{The matrix-valued Crouzeix problem}

The double-layer map is completely positive, so its amplification to
matrix-valued functions preserves positivity without loss.  The new
obstruction is algebraic: scalar diagonal nuisance terms become
noncommuting matrix blocks.  A block, bimodule, or tangential version
of the annihilator construction would therefore be a major advance.
For the complete spectral-set formulation, see the paper of
Davidson, Paulsen, and Woerdeman,
\href{https://arxiv.org/abs/1612.05683}{arXiv:1612.05683}.

\subsection{The abstract Crouzeix conjecture}

Clou\^atre, Ostermann, and Ransford conjecture that additional unital
antilinear structure should force the constant $2$; see
\href{https://arxiv.org/abs/2011.10422}{arXiv:2011.10422}.
The present mechanism establishes this type of conclusion when the
adjoint correction lies in a semisimple annihilable module.  A broader
criterion guaranteeing an annihilator could connect the two methods.

\subsection{Scaled \texorpdfstring{$q$}{q}-numerical ranges}

Recent work extends the Crouzeix--Palencia framework to scaled
$q$-numerical ranges and proposes a corresponding Crouzeix-type
conjecture; see
\href{https://arxiv.org/abs/2603.15536}
{O'Loughlin--Rani, arXiv:2603.15536}.
The associated kernel has a comparatively simple nuisance term, so
one can search for $q$-dependent pairs of Gramian weights and then
optimize their scale as in \eqref{eq:optimal-radius}.

\subsection{Annular and multiply connected spectral constants}

Annulus theory combines double-layer positivity with coupled
information about $T$ and $T^{-1}$; see
\href{https://arxiv.org/abs/2307.13387}
{Jury--Tsikalas, arXiv:2307.13387}.
Current work on quantum-annulus constants obtains asymptotically sharp
bounds; see
\href{https://arxiv.org/abs/2505.06230}
{Pascoe, arXiv:2505.06230}.
This setting naturally suggests paired forward and backward Gramians,
together with more than one cancellation anchor.

\subsection{Multivariable Schur--Agler problems}

Positive-kernel descriptions of the Schur--Agler norm involve
diagonalizable tuples and semidefinite optimization; see
\href{https://arxiv.org/abs/2602.14047}
{Hartz--Wang, arXiv:2602.14047}.
Kernel annihilation might eliminate nonunique Agler components before
a multi-index Gramian argument is applied.  This direction is more
speculative because several positive kernels and several Stein
recursions occur simultaneously.

\section{Why adjoint alignment is indispensable}

The orientation of the nuisance algebra cannot be discarded.  For
$M>0$, let
\begin{equation}
 T_M=\begin{pmatrix}0&M\\0&1/2\end{pmatrix},
 \qquad H(z)\equiv I.
 \label{eq:orientation-counterexample}
\end{equation}
Then
\begin{equation}
 \sigma(T_M)\subset\D,
 \qquad \RePart H=I,
 \label{eq:counterexample-positive}
\end{equation}
and
\begin{equation}
 H(z)-(I-zT_M)^{-1}
 =-\sum_{k\geq1}z^kT_M^k
 \in\alg(T_M).
 \label{eq:wrong-algebra-membership}
\end{equation}
Nevertheless, \(\|T_M\|\to\infty\) as $M\to\infty$.  Thus replacing
\(\alg(T^*)\) by \(\alg(T)\) destroys the conclusion even though both
algebras are commutative.  The adjoint orientation is exactly what
makes the nuisance term compatible with the annihilating test data.

\section{Research program}

The reusable content of the mechanism is therefore the following
five-step program:
\begin{enumerate}
 \item Find a positive kernel.
 \item Identify an adjoint-aligned nuisance module.
 \item Construct test vectors that annihilate the entire module.
 \item Arrange the surviving terms as an ordered difference of
       weighted Gramians.
 \item Verify detectability and a Stein recursion.
\end{enumerate}
The method is credible precisely where all five steps can be carried
out without estimating the nuisance term.  The complete Crouzeix
problem and scaled $q$-numerical-range problems are the most direct
targets; annular and multivariable variants require additional
forward--backward or multi-index structure.

\end{document}
